%~~~~~~~~~~~~~~~~~~~~~~~~~~~~~~~~~~~~~~~~~~~~~~~~~~~~~~~~~~~~~~~~~~~~~~~~~~~~~~~~~~~~~~~~~%
%	                                     ABSTRACT   	                                  %
%~~~~~~~~~~~~~~~~~~~~~~~~~~~~~~~~~~~~~~~~~~~~~~~~~~~~~~~~~~~~~~~~~~~~~~~~~~~~~~~~~~~~~~~~~%
\begin{flushleft} % This section is not essential for the abstract
    \setlength{\parskip}{0pt}
        %\bigskip
    {\centering{{\Large{\bf{ABSTRACT}}}} \par}
    \medskip
    \vspace{6pt}
    \hrule \vspace{0.4cm}
		  Student Name:	 \textbf{\ThesisAuthor}	\hfill Reg. No.: \textbf{\RegNumber} \\
		Degree:  \textbf{\DegreeDisplay} \hfill Dept.:  \textbf{School of \Department} \\
		Thesis Title: \textbf{\ThesisTitle}\\
		Thesis Supervisor:  \textbf{\ProjectSupervisor}\\
		Date of thesis submission: \textbf{{ \SubmissionDate}\\[0.4cm] }
	\hrule \vspace{1.5cm} 
\end{flushleft}
\noindent
Radio frequency interference (RFI) from human technology corrupts radio
astronomy data and must be flagged before analysis. Treating RFI flagging as
image segmentation, this project reproduces the U-Net method of Akeret et al.\
(\tfunet{}), designs an improved network, and tests both on a synthetic
dataset and on real LOFAR observations labelled by a human expert. On synthetic
data \tfunet{} first scored an \Fone{} of 0.39; controlled experiments traced this
to class weighting combined with its ReLU output layer and to per-image
normalisation, and correcting both raised it to 0.93 without changing the
network. On real LOFAR data, the improved network, \hybrid{}, reaches a maximum
\Fone{} of $0.6603 \pm 0.0040$ with 593{,}842 parameters, above the best published
result of 0.5979. However, removing the components it was designed around---strip
convolutions, channel attention and residual blocks---leaves its score
unchanged, and a conditional random field refinement adds nothing. Checking the
public datasets revealed test images duplicated in the training sets of both
HERA and LOFAR, and the same network scores 0.44 lower on real data than on
synthetic data. The results show that synthetic benchmarks and single training
runs can both mislead, and that claims about why a model works need to be
tested directly.

\vspace{4cm}

{\large \textbf{Keywords: }}\par{\Large \ThesisKeywords}

\cleardoublepage